\section{Problem Formulation}
\label{sec:formulation}

\subsection{Fraud Detection as Learning to Defer}
\label{sec:fraud_l2d}

We model insurance fraud detection as a learning-to-defer problem. Each claim $\xx \in \mathcal{X} \subseteq \RR^d$ has a binary label $y \in \{0, 1\}$ (legitimate or fraudulent). The system must choose one of three actions: $a \in \{\text{clear}, \text{flag}, \deferop\}$. When the system defers, a human fraud investigator $M$ examines the claim with accuracy parameterised by $\alpha \in (0, 1]$.

The \textit{system-level} prediction is:
\begin{equation}
    \hat{y}_{\text{sys}}(\xx) = \begin{cases}
        h(\xx) & \text{if } r(\xx) = 0 \quad (\text{classifier decides}), \\
        M(\xx) & \text{if } r(\xx) = 1 \quad (\text{investigator decides}).
    \end{cases}
\end{equation}
The \textit{coverage} is the fraction of claims decided by the classifier: $\text{Coverage} = 1 - \frac{1}{N}\sum_{i} r(\xx_i)$.

\subsection{Two-Stage Risk-Sensitive Deferral}
\label{sec:risk_objective}

Since the base classifier $h$ is typically pre-trained or provided by a third party, we decouple classifier training from deferral optimisation: Stage~1 trains $h$ via standard supervised learning, and Stage~2 selects the rejector $r$ by optimising a CVaR-penalised objective over the deferral threshold. Formally:

\begin{definition}[CVaR-Guided Deferral Objective]
\label{def:cvar_l2d}
Given a fixed classifier $h$, risk level $\delta \in (0,1]$, and trade-off weight $\lambda > 0$, the deferral threshold $\tau$ is selected by:
\begin{equation}
\begin{aligned}
\tau^*
&= \argmax_{\tau} \;
\underbrace{\mathrm{Acc}_{\mathrm{sys}}(h, r_\tau)}_{\text{system accuracy}} \\
&\quad - \lambda \underbrace{\cvar_\delta\bigl[\ell(\xx, y, h(\xx)) \cdot \mathbf{1}\{r_\tau(\xx) = 0\}\bigr]}_{\text{tail risk on retained claims}}
\end{aligned}
\label{eq:cvar_l2d}
\end{equation}

where $r_\tau(\xx) = \mathbf{1}[U(\xx) > \tau]$ is a threshold rejector based on an uncertainty signal $U(\xx)$ (defined in \Cref{sec:method_bayes}), and $\ell(\xx, y, h(\xx)) = \mathbf{1}[h(\xx) \neq y]$ is the 0-1 loss on non-deferred instances.
\end{definition}

The first term rewards effective deferral of hard cases; the second penalises configurations where retained predictions carry high worst-case error. The CVaR term admits the Rockafellar--Uryasev dual \citep{rockafellar2000optimization}:
\begin{equation}
    \cvar_\delta[Z] = \inf_{\eta \in \RR} \left\{ \eta + \frac{1}{\delta}\EE\bigl[(Z - \eta)^+\bigr] \right\},
    \label{eq:cvar_dual}
\end{equation}
which is convex in $\eta$. Although the one-dimensional threshold sweep in \Cref{eq:cvar_l2d} does not directly exploit this convexity, the coherence and smoothness of CVaR make the objective well-behaved compared to, e.g., VaR.

\begin{proposition}[Tail-Risk Reduction via Deferral]
\label{prop:cvar_bound}
Let $r$ be any rejector that defers at least one sample. Denote $\mathcal{S}_r = \{i : r(\xx_i) = 0\}$ the retained set. If every deferred sample has loss $\ell_i \geq \var_\delta(\ell)$ (i.e., deferral targets the upper tail), then:
\begin{equation}
    \cvar_\delta\bigl[\ell_i \mid i \in \mathcal{S}_r\bigr] \leq \cvar_\delta\bigl[\ell_i\bigr].
\end{equation}
\end{proposition}
The proof is given in \Cref{app:math}. For 0-1 loss, the condition simplifies: since $\ell_i \in \{0,1\}$, the proposition requires that deferred samples include a sufficient proportion of misclassified instances. We verify this empirically in our experiments, the Spearman rank correlation between posterior entropy $H(\pi_F)$ and the 0-1 loss is $\rho > 0.3$ ($p < 0.001$), confirming that the uncertainty signal preferentially identifies misclassified claims (\Cref{sec:exp_results}).

\subsection{Connection to Constrained Optimisation}
\label{sec:cmdp}

The deferral objective in \Cref{eq:cvar_l2d} can be viewed as the Lagrangian relaxation of a constrained problem. Define:
\begin{equation}
\begin{aligned}
\max_{r} \;& \mathrm{Acc}_{\mathrm{sys}}(h, r) \\
\text{s.t.} \;& \cvar_\delta\bigl[\ell(\xx, y, h(\xx)) \cdot \mathbf{1}\{r(\xx) = 0\}\bigr] \leq \epsilon
\end{aligned}
\label{eq:cmdp}
\end{equation}

Under Slater's condition (which holds when there exists a feasible $r$ achieving $\cvar_\delta < \epsilon$, e.g., by deferring all claims), strong duality gives $\max_r \min_{\lambda \geq 0} \{\mathrm{Acc}_{\mathrm{sys}} + \lambda(\epsilon - \cvar_\delta)\} = \min_{\lambda \geq 0} \max_r \{\mathrm{Acc}_{\mathrm{sys}} + \lambda(\epsilon - \cvar_\delta)\}$, and the optimal $\lambda^*$ recovers the constraint boundary. Thus, \Cref{eq:cvar_l2d} with a specific $\lambda$ corresponds to a particular risk tolerance $\epsilon$, connecting our framework to constrained optimisation \citep{altman1999constrained}. In the sequential setting where claims arrive over time and investigator capacity is finite, this extends to a constrained MDP, motivating online deferral with budget constraints as future work.
